\documentclass[10pt,a4paper]{article}
\usepackage[a4paper,margin=20mm]{geometry}
\usepackage[T1]{fontenc}
\usepackage[utf8]{inputenc}
\usepackage[british]{babel}
\usepackage{amsmath,amssymb}
\usepackage{newtxtext,newtxmath}
\usepackage{graphicx}
\usepackage{float}
\usepackage{booktabs,longtable,array,ragged2e}
\usepackage{setspace}
\usepackage{csquotes}
\usepackage{xurl}
\usepackage[hidelinks]{hyperref}
\newcolumntype{L}[1]{>{\RaggedRight\arraybackslash}p{#1}}
\setlength\LTleft{0pt}
\setlength\LTright{0pt plus 1fill}
\setlength{\emergencystretch}{3em}
\setlength{\parindent}{0pt}
\setlength{\parskip}{0.35em}
\newcommand{\vect}[1]{\boldsymbol{#1}}
\renewcommand{\thefigure}{S\arabic{figure}}
\renewcommand{\thetable}{S\arabic{table}}
\begin{document}
\begin{center}
{\Large\bfseries Supplementary material}\par
\vspace{0.4em}
{\large Recovering regional myocardial dysfunction from routine cine CMR: methods, validation
evidence and open questions}\par
\vspace{0.3em}
{Ziying Huo}\par
\end{center}

\section*{S1. Search record}

The review is a critical narrative review with a structured search, not a systematic review. The
search was designed to map method families and to test the review's central question; it was not
designed to estimate a pooled effect or to guarantee exhaustive retrieval. The evidence cutoff was
1~September 2026. Screening was performed once, without independent duplicate human screening.

\textbf{Sources.} PubMed and PubMed Central/Europe PMC; OpenAlex; Crossref; arXiv. These were
supplemented by backward and forward citation expansion from core method, comparison and
standardisation papers, and by direct retrieval of official dataset, repository, guideline and
regulatory records.

\textbf{Themes.} (i) cine CMR motion estimation and deformable registration; (ii) myocardial strain
measurement and its specification; (iii) tagging, DENSE and prescribed-motion references;
(iv) regional and post-infarction strain; (v) accelerated and deep-learning cine reconstruction.
Learning-based methods were searched from 2016 onwards; earlier registration and tagging work was
retained where it remains the reference point for later methods.

\textbf{Query log.} The exact query strings, database caps and hit counts from the structured gap
search are supplied in the machine-readable file \path{supplementary/gap-search-query-log.csv}. The
queries rely on exact phrases such as \enquote{cine CMR} or \enquote{cine MRI} and do not document a
parallel controlled-vocabulary expansion, so the search can establish whether a given conjunction
was retrieved by the named queries but cannot estimate recall or the prevalence of qualifying
studies. The file records the queries so that they can be re-entered by hand; it is not an
executable script, and no screening flow counts are reported.

\textbf{Post-cutoff update.} One study published after the cutoff (Naqizadeh Jahromi et al., Biomech
Model Mechanobiol 2026;25(5):103) was appraised separately as a targeted update and is identified as
such wherever it is used in the main text. It was not included in the original query or screening
counts.

\textbf{Reference verification.} Every entry in the reference list was checked against Crossref,
Europe PMC, DataCite or the arXiv application programming interface for authors, title, year, venue,
volume and pages, and digital object identifier or arXiv identifier. The entries newly added for
this version were verified when they were added; the entries retained from the previous version were
re-verified in a second pass against Crossref, the arXiv application programming interface or
DataCite, which corrected three records. All 112 entries with a registered identifier resolve and
match on authors, title, year, venue, volume and pages, with four documented notes where a publisher
record differs from the entry by design: two preprints or articles cited by the version or issue of
record rather than by the earliest online date, one article number in place of a page range, and one
range completed from Europe PMC where Crossref lists only the first page. Six entries have no
registered identifier because they are regulatory documents, software repositories, a challenge page
or conference proceedings without a digital object identifier; each is cited by its stable URL with
the access date. Records that could not be verified were not cited. Verification confirms that a record
exists and what it is; it does not validate the study it describes.

\textbf{Reading depth.} Each quantitative or comparative statement in the main text is attributed to
a specific source location, and the depth at which that source was read is recorded in the
accompanying claim audit. Sources read at abstract level support abstract-level statements only.

\section*{S2. Separating a training-time change from an inference-time control}

Comparisons that vary a prior must distinguish two different experiments, because they require
different controls (main text, \S6.3). Changing a weight in the \emph{training}
objective changes the trained model and therefore requires matched retraining under controlled
splits and training conditions; changing an exposed \emph{inference-time} setting leaves the trained
model fixed and is only possible where such a control is implemented. Supplementary
Figure~\ref{fig:supp-controls} sets out both, together with the parts of the comparison that must be
held fixed in either case.

\begin{figure}[H]
\centering
\includegraphics[width=\textwidth,keepaspectratio]{figures/p6_FigureS1_training_vs_inference_control.pdf}
\caption{Training-time changes and inference-time controls require different comparisons.
\textbf{a}, changing a training-objective weight requires matched retraining under controlled splits
and training conditions. \textbf{b}, changing an inference-time setting is meaningful only when the
trained model exposes that implemented control; the model otherwise remains fixed. Cases, splits,
measurement definition, evaluation reference, losses and tolerances form the fixed comparison
contract, where \enquote{losses} refers to evaluation losses rather than the training weight varied
in panel a. Model or setting selection must be completed without using the final held-out test set.
This is a proposed comparison design, not an experiment, a performance result or an optimum.}
\label{fig:supp-controls}
\end{figure}

\section*{S3. Reporting a conditional recoverability claim}

This section states the bookkeeping that a conditional recoverability claim requires. It is a
reporting contract, not a validated rule, and no tolerance given here is empirically justified by
this review.

Let $g$ index a regional attribute (amplitude, location, extent or timing), let $L_g$ be the loss
chosen for that attribute, and let $\delta_g$ be a tolerance fixed before the final evaluation. For a
held-out case $i$ with a reference value available, the attribute is \emph{within tolerance} when
$L_g(i)\le\delta_g$. If the method also provides a reporting or abstention rule $r(i)\in\{0,1\}$ that
depends only on observable quality-control signals and not on the unknown true error, then the
quantities that must be reported together are
\[
\text{coverage } = \frac{1}{N}\sum_{i=1}^{N} r(i),
\qquad
\text{tolerance success among reported cases } = \frac{\sum_{i} r(i)\,\mathbb{1}\{L_g(i)\le\delta_g\}}{\sum_{i} r(i)},
\]
together with the false-acceptance rate $\sum_i r(i)\mathbb{1}\{L_g(i)>\delta_g\}/\sum_i r(i)$ and
the coverage within each prespecified subgroup. Reporting the second quantity without the first is
uninformative, because rejecting almost every case drives it towards one; reporting the first
without the second is equally uninformative. This is the standard selective-prediction
decomposition of error conditional on reporting~\cite{geifman2019selectivenet}; it is a reporting
framework, not evidence about CMR measurement accuracy.

A complete claim therefore names: the acquisition and image-quality regime; the abnormality class
and the range of its parameters; the estimator, its version and the setting of any exposed control;
the complete measurement specification (Table~2 of the main text); the evaluation reference and its
own uncertainty; the attribute $g$, the loss $L_g$ and the tolerance $\delta_g$; and the coverage at
which the result holds. Changing any of these changes the claim.

\begin{thebibliography}{1}
\bibitem{geifman2019selectivenet}
Geifman Y, El-Yaniv R. SelectiveNet: a deep neural network with an integrated reject option.
\emph{Proc 36th Int Conf Mach Learn} 2019;97:2151--2159.
\end{thebibliography}

\end{document}
